% =============================================================================
% 4 Entwicklung des transformatorgestützten Stellkonzepts
% =============================================================================
\section{Entwicklung des transformatorgestützten Stellkonzepts}
\label{chap:entwicklung_stellkonzept}

Die Analyse des bestehenden Erwärmungsprüfstands hat gezeigt, dass die
Ströme der einzelnen MCC-Abgänge während der Erwärmungsprüfung nicht
unabhängig voneinander nachgeführt werden können. Die bestehende
Regelung hält lediglich den unmittelbar hinter dem
Hochstromtransformator gemessenen Gesamtstrom auf seinem Sollwert.
Temperaturbedingte Änderungen der einzelnen Lastzweige müssen dagegen
durch ein manuelles Verstellen der V2A-Widerstände ausgeglichen werden.

Ausgehend von den in Abschnitt~\ref{sec:anforderungen_optimierung}
abgeleiteten Anforderungen wird im Folgenden ein transformatorgestütztes
Stellkonzept entwickelt. Ziel des Konzepts ist es, für jede Phase eines
MCC-Abgangs eine separat beeinflussbare Stellgröße bereitzustellen.
Hierdurch soll der jeweilige Modulstrom während der Erwärmungsprüfung
automatisch nachgeführt werden können, ohne die vorhandenen
V2A-Widerstände händisch verstellen zu müssen.

Das Konzept beruht auf einem zusätzlichen Transformator, der im
Folgenden als Trafo~2 bezeichnet wird. Auf dessen Sekundärseite befindet
sich ein motorisch verstellbarer Widerstand
\(R_{\mathrm{Stell}}\). Durch die Impedanztransformation des
Transformators kann der auf der Sekundärseite verwendete
Widerstandsbereich an die wesentlich kleineren Impedanzen des
Hochstromkreises angepasst werden. Damit lassen sich mechanisch
handhabbare Stellwiderstände verwenden, während primärseitig eine
ausreichende Beeinflussung des Modulstroms erreicht wird.

Für die Abgänge mit Bemessungsströmen von
\SI{250}{\ampere}, \SI{400}{\ampere} und
\SI{630}{\ampere} wird Trafo~2 als zusätzlicher Strompfad parallel zum
vorhandenen V2A-Widerstand angeordnet. Bei den Abgängen mit
\SI{40}{\ampere} und \SI{70}{\ampere} entfällt der V2A-Widerstand.
Der Modulstrom wird bei diesen Abgängen vollständig über den
Primärzweig von Trafo~2 geführt und durch den sekundärseitigen
Stellwiderstand beeinflusst.

Da sämtliche zwölf MCC-Abgänge dreiphasig ausgeführt sind und jede
Phase separat eingestellt werden soll, umfasst das Gesamtkonzept

\begin{equation}
    N_{\mathrm{Stellkanäle}}
    =
    12 \cdot 3
    =
    36
    \label{eq:anzahl_stellkanäle}
\end{equation}

phasenselektive Stellkanäle. Die dreiphasigen MCC-Abgänge bleiben dabei
jeweils eine funktionale Einheit. Die Anzahl von 36 bezeichnet daher
nicht die Anzahl der Abgänge, sondern die Anzahl der einzeln zu
beeinflussenden Phasenströme.

Dieses Kapitel beschreibt zunächst die Grundidee und das elektrische
Wirkprinzip des Konzepts. Anschließend werden die Ausführungen für die
verschiedenen Stromklassen, das Mess- und Regelungskonzept sowie die
erforderliche Dimensionierung und Stellreserve behandelt.


% -----------------------------------------------------------------------------
% 4.1 Grundidee und Wirkprinzip
% -----------------------------------------------------------------------------
\subsection{Grundidee und Wirkprinzip}
\label{sec:grundidee_wirkprinzip}

Die Grundidee des entwickelten Konzepts besteht darin, den Strom eines
MCC-Abgangs nicht mehr ausschließlich durch die mechanische Einstellung
des V2A-Widerstands festzulegen. Stattdessen wird ein zusätzlicher,
elektrisch beeinflussbarer Strompfad geschaffen. Dieser Zusatzpfad
besteht primärseitig aus der Wicklung von Trafo~2. Auf der galvanisch
getrennten Sekundärseite wird der Transformator mit dem motorisch
verstellbaren Widerstand \(R_{\mathrm{Stell}}\) belastet.

Bei den Abgängen mit Bemessungsströmen von
\SI{250}{\ampere}, \SI{400}{\ampere} und
\SI{630}{\ampere} liegt die Primärwicklung von Trafo~2 parallel zum
vorhandenen V2A-Widerstand. Der Modulstrom einer Phase teilt sich damit
auf den V2A-Zweig und den Primärzweig von Trafo~2 auf. Für den Abgang
\(k\) und die Phase \(\varphi\) gilt

\begin{equation}
    I_{k,\varphi}
    =
    I_{\mathrm{V2A},k,\varphi}
    +
    I_{\mathrm{T2,1},k,\varphi},
    \qquad
    \varphi \in \{\mathrm{L1},\mathrm{L2},\mathrm{L3}\}.
    \label{eq:stromaufteilung_v2a_trafo2}
\end{equation}

Dabei bezeichnet \(I_{k,\varphi}\) den gesamten Phasenstrom des
MCC-Abgangs, \(I_{\mathrm{V2A},k,\varphi}\) den Strom durch den
V2A-Widerstand und \(I_{\mathrm{T2,1},k,\varphi}\) den Primärstrom von
Trafo~2.

Der sekundärseitige Stellwiderstand wird über das
Übersetzungsverhältnis von Trafo~2 auf die Primärseite transformiert.
Für den idealisierten Transformator ergibt sich die auf die Primärseite
bezogene Stellimpedanz zu

\begin{equation}
    R_{\mathrm{Stell}}'
    =
    \left(
        \frac{N_1}{N_2}
    \right)^2
    R_{\mathrm{Stell}},
    \label{eq:stellwiderstand_transformation}
\end{equation}

wobei \(N_1\) und \(N_2\) die Windungszahlen der Primär- und
Sekundärwicklung von Trafo~2 bezeichnen. Das Übersetzungsverhältnis wird
so gewählt, dass ein auf der Sekundärseite technisch gut realisierbarer
Widerstandsbereich auf die kleinen Impedanzwerte des primärseitigen
Hochstromkreises transformiert wird.

Unter Vernachlässigung der inneren Transformatorimpedanz kann der
Primärstrom des Zusatzpfads näherungsweise durch

\begin{equation}
    I_{\mathrm{T2,1},k,\varphi}
    \approx
    \frac{U_{k,\varphi}}
    {R_{\mathrm{Stell},k,\varphi}'}
    \label{eq:primaerstrom_trafo2_idealisiert}
\end{equation}

beschrieben werden. \(U_{k,\varphi}\) bezeichnet dabei die Spannung
zwischen dem Eingang des jeweiligen Lastzweigs und dessen
Sternpunkt. Gleichung~\eqref{eq:primaerstrom_trafo2_idealisiert} dient
zunächst der Beschreibung des grundlegenden Wirkzusammenhangs. Die
Wicklungswiderstände, Streuinduktivitäten und die Kurzschlussimpedanz
von Trafo~2 werden bei der späteren Modellbildung berücksichtigt.

Wird \(R_{\mathrm{Stell}}\) verkleinert, verringert sich nach
Gleichung~\eqref{eq:stellwiderstand_transformation} auch dessen auf die
Primärseite transformierter Widerstand. Dadurch steigt der Primärstrom
von Trafo~2. Bei den großen Abgängen erhöht sich somit der über den
Zusatzpfad fließende Stromanteil und nach
Gleichung~\eqref{eq:stromaufteilung_v2a_trafo2} der gesamte
Modulstrom. Umgekehrt führt eine Vergrößerung von
\(R_{\mathrm{Stell}}\) zu einer Verringerung des Zusatzstroms.

Der V2A-Widerstand wird während dieser Nachführung nicht verstellt. Er
bildet bei den großen Abgängen weiterhin den Grundlastpfad, während
Trafo~2 und \(R_{\mathrm{Stell}}\) den regelbaren Anteil des
Modulstroms bereitstellen. Damit muss der Zusatzpfad nicht den gesamten
Bemessungsstrom des Abgangs führen, sondern lediglich den zur
Einstellung und zum Ausgleich der thermischen Stromdrift erforderlichen
Stromanteil.

Bei den Abgängen mit Bemessungsströmen von
\SI{40}{\ampere} und \SI{70}{\ampere} wird kein V2A-Widerstand
verwendet. Der Phasenstrom fließt vollständig durch die Primärwicklung
von Trafo~2. Für diese Ausführung gilt daher

\begin{equation}
    I_{k,\varphi}
    =
    I_{\mathrm{T2,1},k,\varphi}.
    \label{eq:strom_kleinmodule_trafo2}
\end{equation}

Die Einstellung des vollständigen Modulstroms erfolgt in diesem Fall
über den sekundärseitigen Stellwiderstand. Hieraus ergeben sich andere
Anforderungen an den Stellbereich und die Belastbarkeit von Trafo~2
und \(R_{\mathrm{Stell}}\) als bei den parallel zum V2A-Widerstand
betriebenen Zusatzpfaden. Die beiden Ausführungsvarianten werden daher
in den folgenden Abschnitten getrennt betrachtet.

Die Verstellung von \(R_{\mathrm{Stell}}\) erfolgt motorisch. Dadurch
steht für jeden Phasenstrom eine automatisierbare Stellgröße zur
Verfügung. Zusammen mit der phasenweisen Strommessung kann somit ein
geschlossener Regelkreis aufgebaut werden, in dem der jeweilige
Modulstrom die Regelgröße und die Position des Schleifkontakts von
\(R_{\mathrm{Stell}}\) die Stellgröße darstellt. Die vorhandene
übergeordnete Regelung des Gesamtstroms bleibt weiterhin wirksam.
Dadurch entsteht ein gekoppeltes Mehrgrößensystem aus der zentralen
Gesamtstromregelung und den 36 unterlagerten phasenselektiven
Stellkanälen.



% -----------------------------------------------------------------------------
% 4.2 Zusatzpfad aus Trafo 2 und Stellwiderstand
% -----------------------------------------------------------------------------
\subsection{Zusatzpfad aus Trafo~2 und Stellwiderstand}
\label{sec:zusatzpfad_trafo2}

Der regelbare Zusatzpfad besteht aus der Primärwicklung von Trafo~2 und
dem auf dessen Sekundärseite angeschlossenen Stellwiderstand
\(R_{\mathrm{Stell}}\). Zwischen dem primärseitigen Hochstromkreis und
dem sekundärseitigen Stellkreis besteht keine galvanische Verbindung.
Die Beeinflussung des Primärstroms erfolgt ausschließlich über die
magnetische Kopplung der beiden Transformatorwicklungen.

Die Aufgabe von Trafo~2 besteht nicht darin, dem MCC-Abgang Energie aus
einer zusätzlichen Spannungsquelle zuzuführen. Die für den Zusatzstrom
erforderliche Leistung wird weiterhin durch den bestehenden
Erwärmungsprüfstand bereitgestellt. Trafo~2 dient vielmehr dazu, die
sekundärseitige Lastimpedanz auf ein für den primärseitigen
Hochstromkreis geeignetes Niveau zu transformieren. Dadurch kann ein
mechanisch und elektrisch handhabbarer Stellwiderstand verwendet
werden, obwohl die Impedanzen im primärseitigen Hochstromkreis im
Milliohmbereich liegen.

Das Übersetzungsverhältnis von Trafo~2 wird definiert als

\begin{equation}
    u_{\mathrm{T2}}
    =
    \frac{N_1}{N_2}
    =
    \frac{U_1}{U_2}
    =
    \frac{I_2}{I_1},
    \label{eq:uebersetzungsverhaeltnis_trafo2}
\end{equation}

wobei \(N_1\) und \(N_2\) die Windungszahlen, \(U_1\) und \(U_2\) die
Effektivwerte der Spannungen sowie \(I_1\) und \(I_2\) die
Effektivwerte der Ströme auf der Primär- und Sekundärseite bezeichnen.
Die Beziehungen gelten zunächst für den idealisierten Transformator.
Die grundlegenden Transformatorbeziehungen und die
Impedanztransformation wurden bereits in
Abschnitt~\ref{sec:transformator_impedanztransformation} behandelt.

Der auf der Sekundärseite angeschlossene Stellwiderstand erscheint auf
der Primärseite mit dem transformierten Wert

\begin{equation}
    R_{\mathrm{Stell}}'
    =
    u_{\mathrm{T2}}^2
    R_{\mathrm{Stell}}.
    \label{eq:rstell_auf_primaerseite}
\end{equation}

Für eine Spannungserhöhung von der Primär- zur Sekundärseite gilt
\(N_2>N_1\) und damit \(u_{\mathrm{T2}}<1\). Der sekundärseitige
Stellwiderstand wird folglich mit dem Faktor
\(u_{\mathrm{T2}}^2\) auf die Primärseite transformiert. Ein
vergleichsweise großer und technisch gut einstellbarer
Sekundärwiderstand entspricht dadurch einem wesentlich kleineren
Widerstand im Hochstromkreis.

Der reale Transformator weist zusätzlich Wicklungswiderstände,
Streuinduktivitäten und einen Magnetisierungsstrom auf. Für die
Beschreibung des belasteten Zusatzpfads kann die auf die Primärseite
bezogene Gesamtimpedanz vereinfacht als

\begin{equation}
    \underline{Z}_{\mathrm{Zus}}
    =
    \underline{Z}_{\mathrm{T2}}
    +
    u_{\mathrm{T2}}^2 R_{\mathrm{Stell}}
    \label{eq:gesamtimpedanz_zusatzpfad}
\end{equation}

angegeben werden. Darin fasst
\(\underline{Z}_{\mathrm{T2}}\) die für das Betriebsverhalten
maßgebenden inneren Transformatorimpedanzen zusammen. Der Unterstrich
kennzeichnet eine komplexe Wechselstromgröße. Der zugehörige
Primärstrom des Zusatzpfads ergibt sich zu

\begin{equation}
    \underline{I}_{\mathrm{T2,1}}
    =
    \frac{\underline{U}_{1}}
    {\underline{Z}_{\mathrm{T2}}
    + u_{\mathrm{T2}}^2 R_{\mathrm{Stell}}}.
    \label{eq:primaerstrom_zusatzpfad}
\end{equation}

Gleichung~\eqref{eq:primaerstrom_zusatzpfad} zeigt den grundlegenden
Stellzusammenhang. Durch eine Verringerung von
\(R_{\mathrm{Stell}}\) sinkt die auf die Primärseite transformierte
Gesamtimpedanz des Zusatzpfads. Bei näherungsweise konstanter
Primärspannung steigt dadurch der Primärstrom von Trafo~2. Eine
Vergrößerung von \(R_{\mathrm{Stell}}\) bewirkt entsprechend eine
Verringerung des Primärstroms.

Die Stellrichtung lässt sich damit qualitativ durch

\begin{equation}
    R_{\mathrm{Stell}} \downarrow
    \quad\Longrightarrow\quad
    R_{\mathrm{Stell}}' \downarrow
    \quad\Longrightarrow\quad
    I_{\mathrm{T2,1}} \uparrow
    \label{eq:stellrichtung_trafo2}
\end{equation}

beschreiben. Diese monotone Beziehung ist für die spätere
Stromregelung von wesentlicher Bedeutung, da jeder Stellung des
Schleifkontakts innerhalb des nutzbaren Stellbereichs ein bestimmter
sekundärseitiger Widerstand und damit ein bestimmter primärseitiger
Strom zugeordnet werden kann.

Der Stellwiderstand besteht aus einem Widerstandselement mit einem
motorisch verfahrbaren Kontakt. Durch die Veränderung der
Kontaktposition \(x_{\mathrm{Stell}}\) wird die wirksame Länge des
Widerstandsmaterials verändert. Für ein Widerstandselement mit
annähernd konstantem Querschnitt kann der Zusammenhang idealisiert
durch

\begin{equation}
    R_{\mathrm{Stell}}
    \left(x_{\mathrm{Stell}}\right)
    =
    \rho_{\mathrm{Stell}}
    \frac{
        l_{\mathrm{wirksam}}
        \left(x_{\mathrm{Stell}}\right)
    }
    {A_{\mathrm{Stell}}}
    \label{eq:rstell_position}
\end{equation}

beschrieben werden. Dabei sind
\(\rho_{\mathrm{Stell}}\) der spezifische elektrische Widerstand,
\(A_{\mathrm{Stell}}\) der Leiterquerschnitt und
\(l_{\mathrm{wirksam}}\) die von der Kontaktposition abhängige wirksame
Leiterlänge. Bei gleichmäßiger Geometrie ergibt sich näherungsweise ein
linearer Zusammenhang zwischen Kontaktposition und Widerstand. Im
realen Aufbau können Kontaktwiderstände, Anschlusswiderstände und
geometrische Übergangsbereiche zu Abweichungen von diesem idealisierten
Verlauf führen.

Die sekundärseitige Belastung des Stellwiderstands ergibt sich aus

\begin{equation}
    P_{\mathrm{Stell}}
    =
    I_2^2 R_{\mathrm{Stell}}.
    \label{eq:verlustleistung_rstell}
\end{equation}

Obwohl der Sekundärstrom aufgrund der gewählten
Spannungstransformation kleiner als der Primärstrom ist, muss
\(R_{\mathrm{Stell}}\) für die auftretende Verlustleistung und den
Dauerbetrieb während einer Erwärmungsprüfung ausgelegt werden.
Zusätzlich sind die Belastbarkeit von Trafo~2, dessen
Kurzschlussimpedanz sowie die thermischen Grenzen beider Komponenten
zu berücksichtigen. Die quantitative Dimensionierung erfolgt in
Abschnitt~\ref{sec:dimensionierung_stellreserve}.

Der mögliche Primärstrom wird durch die beiden Endstellungen des
Stellwiderstands begrenzt. Mit dem minimalen und maximalen
Stellwiderstand

\begin{equation}
    R_{\mathrm{Stell,min}}
    \leq
    R_{\mathrm{Stell}}
    \leq
    R_{\mathrm{Stell,max}}
    \label{eq:stellbereich_rstell}
\end{equation}

ergibt sich ein zugehöriger Stromstellbereich

\begin{equation}
    I_{\mathrm{T2,1,min}}
    \leq
    I_{\mathrm{T2,1}}
    \leq
    I_{\mathrm{T2,1,max}}.
    \label{eq:stromstellbereich_trafo2}
\end{equation}

Der minimale Widerstand darf nicht zu einem unzulässig hohen
Primär- oder Sekundärstrom führen. Der maximale Widerstand muss
dagegen groß genug sein, um den Zusatzstrom auf den für die
Grundeinstellung erforderlichen Wert zu begrenzen. Ein vollständig
stromloser Zusatzpfad wird durch einen endlichen
\(R_{\mathrm{Stell,max}}\) nicht erreicht, da neben dem
lastabhängigen Strom auch der Magnetisierungsstrom von Trafo~2
verbleibt. Falls eine vollständige Trennung erforderlich ist, müsste
hierfür ein zusätzliches Schaltgerät vorgesehen werden.

Bei den Abgängen mit V2A-Widerstand stellt dieser den überwiegenden
Grundlastpfad bereit. Der Zusatzpfad muss daher nur den für die
Feineinstellung und die Kompensation des thermisch bedingten
Stromrückgangs erforderlichen Stromanteil übernehmen. Bei den
\SI{40}{\ampere}- und \SI{70}{\ampere}-Abgängen fließt dagegen der
vollständige Modulstrom über die Primärwicklung von Trafo~2. Daraus
ergeben sich unterschiedliche Anforderungen an Übersetzungsverhältnis,
Stellbereich und Bemessungsleistung. Die konstruktive Ausführung für
beide Anwendungsfälle wird in den folgenden Abschnitten getrennt
untersucht.




% -----------------------------------------------------------------------------
% 4.3 Ausführung für die Abgänge mit hohen Bemessungsströmen
% -----------------------------------------------------------------------------
\subsection{Ausführung für die Abgänge mit hohen Bemessungsströmen}
\label{sec:ausfuehrung_hohe_modulstroeme}

Bei den Abgängen mit Bemessungsströmen von
\SI{250}{\ampere}, \SI{400}{\ampere} und
\SI{630}{\ampere} bleiben die vorhandenen V2A-Widerstände als
Grundlastpfade erhalten. Trafo~2 bildet für jede Phase einen
zusätzlichen, parallel zum jeweiligen V2A-Widerstand angeordneten
Strompfad. Der Zusatzpfad übernimmt damit nicht den vollständigen
Modulstrom, sondern lediglich den zur Einstellung und Nachführung
erforderlichen Stromanteil.

Die Aufteilung des Stroms erfolgt hinter dem zu prüfenden
MCC-Modul. Dadurch durchfließt der gesamte Phasenstrom zunächst die
inneren Strombahnen, Schaltgeräte und Kontaktstellen des Moduls. Erst
am Ausgang des Moduls teilt sich der Strom auf den V2A-Grundlastpfad
und die Primärwicklung von Trafo~2 auf. Beide Teilströme werden
anschließend am gemeinsamen Sternpunkt wieder zusammengeführt. Somit
entspricht die Summe der Teilströme dem thermisch wirksamen Strom durch
das zu prüfende Modul.

Für den Abgang \(k\) und die Phase \(\varphi\) gilt entsprechend

\begin{equation}
    I_{k,\varphi}
    =
    I_{\mathrm{V2A},k,\varphi}
    +
    I_{\mathrm{T2,1},k,\varphi}.
    \label{eq:strombilanz_grosse_abgaenge}
\end{equation}

Der Index \(k\) bezeichnet einen der Abgänge mit
\SI{250}{\ampere}, \SI{400}{\ampere} oder
\SI{630}{\ampere}, während

\begin{equation}
    \varphi
    \in
    \{\mathrm{L1},\mathrm{L2},\mathrm{L3}\}
    \label{eq:phasenmenge_grosse_abgaenge}
\end{equation}

die betrachtete Phase kennzeichnet. Da drei dreiphasige Abgänge dieser
Stromklasse vorgesehen sind, werden für ihre individuelle Nachführung

\begin{equation}
    N_{\mathrm{Stellkanäle,hoch}}
    =
    3 \cdot 3
    =
    9
    \label{eq:anzahl_stellkanaele_hohe_stroeme}
\end{equation}

phasenselektive Stellkanäle benötigt.

Der Strom durch den V2A-Grundlastpfad wird durch die anliegende
Phasenspannung und den temperaturabhängigen Widerstand bestimmt:

\begin{equation}
    I_{\mathrm{V2A},k,\varphi}
    =
    \frac{U_{k,\varphi}}
    {R_{\mathrm{V2A},k,\varphi}(\vartheta)}.
    \label{eq:v2a_strom_grosse_abgaenge}
\end{equation}

Mit zunehmender Temperatur steigt der elektrische Widerstand des
V2A-Materials. Bei näherungsweise konstanter Spannung nimmt der Strom
durch den V2A-Zweig daher ab. Der hierdurch entstehende Fehlbetrag wird
durch eine Erhöhung des Primärstroms von Trafo~2 ausgeglichen.

Der zur Einhaltung des Bemessungsstroms erforderliche Sollwert des
Zusatzstroms ergibt sich aus der Differenz zwischen dem vorgesehenen
Modulstrom und dem aktuell durch den V2A-Zweig fließenden Strom:

\begin{equation}
    I_{\mathrm{T2,1,soll},k,\varphi}
    =
    I_{\mathrm{N},k}
    -
    I_{\mathrm{V2A},k,\varphi}.
    \label{eq:sollstrom_zusatzpfad_grosse_abgaenge}
\end{equation}

Dabei bezeichnet \(I_{\mathrm{N},k}\) den für den jeweiligen Abgang
vorgesehenen Prüfstrom. Durch das motorische Verkleinern von
\(R_{\mathrm{Stell},k,\varphi}\) wird die auf die Primärseite
transformierte Impedanz des Zusatzpfads reduziert. Infolgedessen steigt
der Primärstrom von Trafo~2, bis die Summe aus V2A-Strom und
Zusatzstrom wieder dem vorgegebenen Modulstrom entspricht.

Die Nachführung lässt sich damit durch die Wirkungskette

\begin{equation}
    \vartheta_{\mathrm{V2A}} \uparrow
    \quad\Longrightarrow\quad
    R_{\mathrm{V2A}} \uparrow
    \quad\Longrightarrow\quad
    I_{\mathrm{V2A}} \downarrow
    \quad\Longrightarrow\quad
    R_{\mathrm{Stell}} \downarrow
    \quad\Longrightarrow\quad
    I_{\mathrm{T2,1}} \uparrow
    \label{eq:wirkungskette_nachfuehrung_grosse_abgaenge}
\end{equation}

darstellen. Trafo~2 und \(R_{\mathrm{Stell}}\) kompensieren damit den
Stromrückgang des V2A-Grundlastpfads, ohne dass dessen mechanische
Kontaktbrücke während der Prüfung nachgestellt werden muss.

Das Konzept kann den Modulstrom ausschließlich durch einen positiven
Zusatzstrom ergänzen. Der Zusatzpfad kann den Strom des
V2A-Grundlastpfads nicht aktiv verringern. Die Grundeinstellung des
V2A-Widerstands muss deshalb so gewählt werden, dass dessen Strom auch
im kalten Zustand unterhalb des vorgesehenen Modulstroms liegt. Unter
Berücksichtigung des minimal erreichbaren Zusatzstroms muss gelten:

\begin{equation}
    I_{\mathrm{V2A,kalt},k,\varphi}
    +
    I_{\mathrm{T2,1,min},k,\varphi}
    \leq
    I_{\mathrm{N},k}.
    \label{eq:bedingung_kaltzustand_grosse_abgaenge}
\end{equation}

Diese Bedingung verhindert, dass der Modulstrom bereits zu Beginn der
Prüfung oberhalb seines Sollwerts liegt, obwohl sich der Zusatzpfad in
seiner Stellung mit dem kleinsten erreichbaren Primärstrom befindet.

Im erwärmten Zustand muss Trafo~2 dagegen einen ausreichend großen
Zusatzstrom bereitstellen können. Für den ungünstigsten vorgesehenen
thermischen Betriebszustand gilt daher:

\begin{equation}
    I_{\mathrm{V2A,warm},k,\varphi}
    +
    I_{\mathrm{T2,1,max},k,\varphi}
    \geq
    I_{\mathrm{N},k}.
    \label{eq:bedingung_warmzustand_grosse_abgaenge}
\end{equation}

Die Bedingungen
\eqref{eq:bedingung_kaltzustand_grosse_abgaenge} und
\eqref{eq:bedingung_warmzustand_grosse_abgaenge} legen den mindestens
erforderlichen Stromstellbereich des Zusatzpfads fest. Zusätzlich ist
eine Stellreserve erforderlich, um Bauteiltoleranzen,
Kontaktwiderstände, Abweichungen der Versorgungsspannung und
unterschiedliche Erwärmungsverläufe der drei Phasen ausgleichen zu
können. Die quantitative Bestimmung dieser Reserve erfolgt in
Abschnitt~\ref{sec:dimensionierung_stellreserve}.

Da der Zusatzpfad parallel zum V2A-Widerstand liegt, liegt an beiden
Zweigen dieselbe Spannung an. Für die idealisierte Parallelschaltung
gilt

\begin{equation}
    U_{\mathrm{V2A},k,\varphi}
    =
    U_{\mathrm{T2,1},k,\varphi}
    =
    U_{k,\varphi}.
    \label{eq:spannung_parallelzweige_grosse_abgaenge}
\end{equation}

Eine Veränderung des Zusatzstroms kann aufgrund der endlichen
Innenimpedanz der gemeinsamen Einspeisung jedoch auch die an den
übrigen Abgängen anliegende Spannung geringfügig beeinflussen. Die
einzelnen Stellkanäle sind daher nicht vollständig voneinander
entkoppelt. Diese gegenseitige Beeinflussung wird bei der Entwicklung
des Mess- und Regelungskonzepts in
Abschnitt~\ref{sec:mess_regelungskonzept} berücksichtigt.

Für die drei Phasen eines Abgangs werden jeweils ein eigener
Transformator beziehungsweise ein eigener magnetisch unabhängiger
Transformatorzweig und ein eigener Stellwiderstand benötigt. Dadurch
können Unterschiede zwischen den Phasen, beispielsweise infolge
abweichender Kontaktwiderstände oder unterschiedlicher thermischer
Bedingungen, separat ausgeglichen werden.

Die Stellwiderstände und die verlustbehafteten Komponenten des
Zusatzpfads sind so anzuordnen, dass ihre Verlustwärme die
Temperaturmessung am Prüfling nicht unzulässig beeinflusst. Insbesondere
die sekundärseitigen Stellwiderstände dürfen keine zusätzliche
Erwärmung des MCC-Feldes verursachen, die nicht aus der elektrischen
Belastung des Prüflings selbst resultiert. Andernfalls könnte das
Ergebnis der Erwärmungsprüfung verfälscht werden.

Die V2A-Widerstände übernehmen bei den großen Abgängen weiterhin den
überwiegenden Anteil des Prüfstroms. Hierdurch werden die erforderliche
Bemessungsleistung von Trafo~2 und die Verlustleistung des
sekundärseitigen Stellwiderstands gegenüber einer vollständigen Führung
des Modulstroms über den Transformator reduziert. Gleichzeitig bleibt
durch den regelbaren Zusatzpfad die Möglichkeit erhalten, den
thermisch verursachten Stromrückgang während der gesamten Prüfdauer
auszugleichen.




% -----------------------------------------------------------------------------
% 4.4 Ausführung für Abgänge mit kleinen Bemessungsströmen
% -----------------------------------------------------------------------------
\subsection{Ausführung für Abgänge mit kleinen Bemessungsströmen}
\label{sec:ausfuehrung_kleine_stroeme}

Die Abgänge mit Bemessungsströmen von \SI{40}{\ampere} und
\SI{70}{\ampere} werden ohne zusätzlichen V2A-Lastwiderstand
ausgeführt. Im Gegensatz zu den in
Abschnitt~\ref{sec:ausfuehrung_hohe_stroeme} betrachteten Abgängen
bildet Trafo~2 somit keinen parallelen Zusatzpfad. Der gesamte
Phasenstrom des jeweiligen MCC-Abgangs fließt durch dessen
Primärwicklung. Die Sekundärwicklung wird wiederum mit dem motorisch
verstellbaren Widerstand \(R_{\mathrm{Stell}}\) belastet.

Für einen Abgang \(k\) und eine Phase \(\varphi\) gilt daher

\begin{equation}
    \underline{I}_{k,\varphi}
    =
    \underline{I}_{\mathrm{T2,1},k,\varphi}.
    \label{eq:strom_kleinabgang}
\end{equation}

Unter Berücksichtigung der auf die Primärseite bezogenen
Transformatorimpedanz kann der Phasenstrom entsprechend
Abschnitt~\ref{sec:zusatzpfad_trafo2} näherungsweise durch

\begin{equation}
    \underline{I}_{k,\varphi}
    =
    \frac{\underline{U}_{k,\varphi}}
    {
        \underline{Z}_{\mathrm{T2},k,\varphi}
        +
        u_{\mathrm{T2},k}^{\,2}
        R_{\mathrm{Stell},k,\varphi}
    }
    \label{eq:strom_kleinabgang_rstell}
\end{equation}

beschrieben werden. Darin bezeichnet
\(\underline{Z}_{\mathrm{T2},k,\varphi}\) die auf die Primärseite
bezogene innere Impedanz von Trafo~2. Das Übersetzungsverhältnis
\(u_{\mathrm{T2},k}\) ist entsprechend
Gleichung~\eqref{eq:uebersetzungsverhaeltnis_trafo2} definiert.

Die Verstellung des Abgangsstroms erfolgt ausschließlich über
\(R_{\mathrm{Stell},k,\varphi}\). Wird der Stellwiderstand verkleinert,
verringert sich die auf die Primärseite transformierte Lastimpedanz.
Dadurch steigt der Primärstrom und somit unmittelbar der Strom des
MCC-Abgangs. Eine Vergrößerung des Stellwiderstands bewirkt den
entgegengesetzten Vorgang. Die Wirkungsrichtung lässt sich damit als

\begin{equation}
    R_{\mathrm{Stell},k,\varphi} \downarrow
    \quad\Longrightarrow\quad
    \left|
        \underline{I}_{k,\varphi}
    \right| \uparrow
\end{equation}

beziehungsweise

\begin{equation}
    R_{\mathrm{Stell},k,\varphi} \uparrow
    \quad\Longrightarrow\quad
    \left|
        \underline{I}_{k,\varphi}
    \right| \downarrow
\end{equation}

angeben.

Während Trafo~2 bei den Abgängen mit hohen Bemessungsströmen lediglich
den durch den V2A-Widerstand geführten Grundstrom ergänzt, muss er bei
den \SI{40}{\ampere}- und \SI{70}{\ampere}-Abgängen den vollständigen
Abgangsstrom führen. Der erforderliche Stellbereich ist deshalb nicht
nur für die Kompensation einer thermisch verursachten Stromänderung
auszulegen. Vielmehr muss der vorgesehene Bemessungsstrom vollständig
innerhalb des durch \(R_{\mathrm{Stell}}\) einstellbaren Strombereichs
liegen.

Für den kleinsten und den größten einstellbaren Strom gilt

\begin{align}
    I_{k,\varphi,\min}
    &=
    \left|
        \frac{\underline{U}_{k,\varphi}}
        {
            \underline{Z}_{\mathrm{T2},k,\varphi}
            +
            u_{\mathrm{T2},k}^{\,2}
            R_{\mathrm{Stell},k,\varphi,\max}
        }
    \right|,
    \label{eq:strom_kleinabgang_min}
    \\[1ex]
    I_{k,\varphi,\max}
    &=
    \left|
        \frac{\underline{U}_{k,\varphi}}
        {
            \underline{Z}_{\mathrm{T2},k,\varphi}
            +
            u_{\mathrm{T2},k}^{\,2}
            R_{\mathrm{Stell},k,\varphi,\min}
        }
    \right|.
    \label{eq:strom_kleinabgang_max}
\end{align}

Damit der jeweilige Bemessungsstrom eingestellt und während der
Erwärmungsprüfung nachgeführt werden kann, muss die Bedingung

\begin{equation}
    I_{k,\varphi,\min}
    <
    I_{\mathrm{N},k}
    <
    I_{k,\varphi,\max}
    \label{eq:stellbedingung_kleine_abgaenge}
\end{equation}

erfüllt sein. Der Bemessungsstrom soll dabei nicht unmittelbar an einer
Grenze des Stellbereichs liegen. Sowohl unterhalb als auch oberhalb des
Arbeitspunkts ist eine Stellreserve vorzusehen, damit
Bauteiltoleranzen, Spannungsabweichungen und temperaturabhängige
Änderungen der inneren Impedanzen ausgeglichen werden können. Die
quantitative Festlegung dieser Reserven erfolgt in
Abschnitt~\ref{sec:dimensionierung_stellreserve}.

Das betrachtete MCC-Feld enthält sechs Abgänge mit einem
Bemessungsstrom von \SI{40}{\ampere} und drei Abgänge mit einem
Bemessungsstrom von \SI{70}{\ampere}. Da jede Phase separat beeinflusst
werden soll, ergeben sich für die \SI{40}{\ampere}-Abgänge

\begin{equation}
    n_{\mathrm{Stell},40}
    =
    6 \cdot 3
    =
    18
    \label{eq:stellkanaele_40a}
\end{equation}

und für die \SI{70}{\ampere}-Abgänge

\begin{equation}
    n_{\mathrm{Stell},70}
    =
    3 \cdot 3
    =
    9
    \label{eq:stellkanaele_70a}
\end{equation}

phasenbezogene Stellkanäle. Insgesamt werden für diese beiden
Modulgruppen somit

\begin{equation}
    n_{\mathrm{Stell,klein}}
    =
    n_{\mathrm{Stell},40}
    +
    n_{\mathrm{Stell},70}
    =
    27
    \label{eq:stellkanaele_kleine_abgaenge}
\end{equation}

Stellkanäle benötigt.

Ob für beide Stromklassen dieselbe Ausführung von Trafo~2 und
\(R_{\mathrm{Stell}}\) verwendet werden kann, ist anhand der
erforderlichen Strombereiche, Übersetzungsverhältnisse und
Verlustleistungen zu überprüfen. Eine identische Ausführung wird daher
nicht vorausgesetzt. Insbesondere müssen die Primärwicklungen von
Trafo~2 für den jeweiligen vollständigen Abgangsstrom im Dauerbetrieb
ausgelegt werden. Gleichzeitig sind die sekundärseitigen
Stellwiderstände hinsichtlich ihres Widerstandsbereichs, ihres
zulässigen Stroms und ihrer thermischen Belastbarkeit zu dimensionieren.

Da bei dieser Ausführung kein V2A-Widerstand vorhanden ist, entfällt
dessen temperaturabhängige Widerstandszunahme als unmittelbare
Störgröße. Temperaturabhängige Veränderungen der Wicklungswiderstände,
der Kontaktstellen und der übrigen Abgangsimpedanzen bleiben jedoch
bestehen. Darüber hinaus können Änderungen der gemeinsamen
Versorgungsspannung den Abgangsstrom beeinflussen. Die phasenbezogene
Stromregelung ist deshalb auch für die kleinen Abgänge erforderlich.
Ihre konzeptionelle Einbindung wird in
Abschnitt~\ref{sec:mess_regelungskonzept} beschrieben.

% -----------------------------------------------------------------------------
% 4.5 Vorgesehenes Mess- und Regelungskonzept
% -----------------------------------------------------------------------------
\subsection{Vorgesehenes Mess- und Regelungskonzept}
\label{sec:mess_regelungskonzept}

Die bisherige Regelung des Erwärmungsprüfstands führt den
Säulenstelltransformator anhand des gemessenen Gesamtstroms nach.
Dieses Regelungsprinzip ist für den bestehenden Prüfaufbau geeignet,
steht bei einer zusätzlichen individuellen Regelung der Modulströme
jedoch in Wechselwirkung mit den neu hinzukommenden Regelkreisen.
Verändert einer der unterlagerten Regler den Widerstand
\(R_{\mathrm{Stell}}\), ändert sich der betreffende Modulstrom und
damit zugleich der von der gemeinsamen Einspeisung bereitgestellte
Gesamtstrom. Eine weiterhin aktive Gesamtstromregelung würde auf diese
Änderung mit einer Anpassung der Versorgungsspannung reagieren. Diese
Spannungsänderung würde wiederum sämtliche parallel angeschlossenen
Abgänge beeinflussen.

Für das entwickelte Optimierungskonzept ist daher vorgesehen, die
bisherige Gesamtstromregelung des Säulenstelltransformators durch eine
Spannungsregelung zu ersetzen. Die vorgelagerte Regelung soll die drei
Phasenspannungen auf der Niederspannungsseite des
Hochstromtransformators auf vorgegebenen Sollwerten halten. Die
individuellen Modulströme werden anschließend durch die
phasenbezogene Verstellung der sekundärseitigen Widerstände
\(R_{\mathrm{Stell}}\) geregelt.

Damit ergibt sich eine zweistufige Regelungsstruktur aus einer
vorgelagerten Spannungsregelung und den nachgelagerten
Stromregelkreisen. Die vorgelagerte Regelung stellt möglichst konstante
elektrische Randbedingungen für die parallel angeschlossenen
Abgangseinheiten bereit. Die nachgelagerten Regelkreise gleichen
dagegen die abgangsspezifischen Abweichungen aus.

Für die vorgelagerte Regelung werden die drei Phasenspannungen zu dem
Vektor

\begin{equation}
    \mathbf{U}
    =
    \begin{bmatrix}
        U_{\mathrm{L1}} &
        U_{\mathrm{L2}} &
        U_{\mathrm{L3}}
    \end{bmatrix}^{\mathrm{T}}
    \label{eq:spannungsvektor}
\end{equation}

zusammengefasst. Die zugehörigen Sollwerte lauten

\begin{equation}
    \mathbf{U}_{\mathrm{soll}}
    =
    \begin{bmatrix}
        U_{\mathrm{L1,soll}} &
        U_{\mathrm{L2,soll}} &
        U_{\mathrm{L3,soll}}
    \end{bmatrix}^{\mathrm{T}}.
    \label{eq:spannungssollwertvektor}
\end{equation}

Die Regelabweichung der vorgelagerten Spannungsregelung ergibt sich
damit zu

\begin{equation}
    \mathbf{e}_{U}
    =
    \mathbf{U}_{\mathrm{soll}}
    -
    \mathbf{U}.
    \label{eq:regelabweichung_spannung}
\end{equation}

Als Stellglieder dienen die motorisch verfahrbaren Stromabnehmer des
Säulenstelltransformators. Da jede Phase über einen eigenen
motorischen Antrieb verfügt, können die Phasenspannungen grundsätzlich
separat beeinflusst werden. Die Spannungsmessung ist für eine spätere
Realisierung am gemeinsamen Einspeisepunkt der Prüfanordnung
vorzusehen. Dadurch wird die tatsächlich an den parallelen
Lastzweigen wirksame Spannung erfasst.

Der Sollwert der Spannungsregelung ist so festzulegen, dass der
vorgesehene Arbeitspunkt der Prüfanordnung erreicht wird und für die
nachgelagerten Stromregelkreise eine ausreichende Stellreserve
verbleibt. Der Gesamtstrom ist nach der Umstellung nicht mehr die
unmittelbare Regelgröße des Säulenstelltransformators. Er muss jedoch
weiterhin messtechnisch überwacht werden, da er für die Belastung des
Leistungsschalterfeldes und der gemeinsamen Strompfade maßgebend ist.

Für jede Phase gilt die Strombilanz

\begin{equation}
    I_{\mathrm{ges},\varphi}
    =
    I_{\mathrm{KB},\varphi}
    +
    \sum_{k=1}^{12} I_{k,\varphi},
    \qquad
    \varphi \in
    \left\{
        \mathrm{L1},
        \mathrm{L2},
        \mathrm{L3}
    \right\}.
    \label{eq:strombilanz_regelungskonzept}
\end{equation}

Dabei bezeichnet \(I_{\mathrm{KB},\varphi}\) den über das
Leistungsschalterfeld und die Kurzschlussbrücke geführten Stromanteil
der jeweiligen Phase. Der Spannungssollwert muss folglich nicht nur
unter Berücksichtigung der MCC-Abgänge, sondern anhand des
vollständigen Prüfaufbaus festgelegt werden.

Die nachgelagerte Stromregelung wird für jeden Abgang und jede Phase
separat vorgesehen. Für den Abgang \(k\) und die Phase \(\varphi\)
lautet die Regelabweichung

\begin{equation}
    e_{I,k,\varphi}
    =
    I_{\mathrm{soll},k,\varphi}
    -
    I_{k,\varphi}.
    \label{eq:regelabweichung_modulstrom}
\end{equation}

Der Sollwert entspricht im regulären Prüfbetrieb dem vorgesehenen
Prüfstrom der jeweiligen Abgangseinheit. Die Stellgröße des
Stromregelkreises ist die Position des motorisch verfahrbaren
Stellwiderstands. Über diese Position wird der Widerstand

\begin{equation}
    R_{\mathrm{Stell},k,\varphi}
    =
    f\left(
        x_{\mathrm{Stell},k,\varphi}
    \right)
    \label{eq:stellwiderstand_position_regelung}
\end{equation}

eingestellt. Darin bezeichnet
\(x_{\mathrm{Stell},k,\varphi}\) die mechanische Position des
zugehörigen Schleifkontakts.

Ist der gemessene Modulstrom kleiner als sein Sollwert, gilt

\begin{equation}
    e_{I,k,\varphi} > 0.
\end{equation}

In diesem Fall muss der wirksame Stellwiderstand verkleinert werden.
Die auf die Primärseite von Trafo~2 transformierte Impedanz nimmt
dadurch ab und der Modulstrom steigt. Ist der gemessene Strom dagegen
größer als der Sollwert, muss \(R_{\mathrm{Stell},k,\varphi}\)
vergrößert werden. Die erforderliche Wirkungsrichtung des
Stromregelkreises lässt sich somit durch

\begin{equation}
    \begin{aligned}
        e_{I,k,\varphi} > 0
        &\quad\Longrightarrow\quad
        R_{\mathrm{Stell},k,\varphi} \downarrow,
        \\[0.5ex]
        e_{I,k,\varphi} < 0
        &\quad\Longrightarrow\quad
        R_{\mathrm{Stell},k,\varphi} \uparrow
    \end{aligned}
    \label{eq:regelrichtung_stellwiderstand}
\end{equation}

beschreiben.

Die Widerstandsverstellung ist auf den physikalisch verfügbaren
Stellbereich zu begrenzen:

\begin{equation}
    R_{\mathrm{Stell},k,\varphi,\min}
    \leq
    R_{\mathrm{Stell},k,\varphi}
    \leq
    R_{\mathrm{Stell},k,\varphi,\max}.
    \label{eq:begrenzung_stellwiderstand_regelung}
\end{equation}

Die Begrenzung verhindert, dass der Regler einen außerhalb des
technisch realisierbaren Bereichs liegenden Widerstandswert anfordert.
Für eine spätere konstruktive Umsetzung sind zusätzlich mechanische
Endlagen, zulässige Verfahrgeschwindigkeiten sowie Strom- und
Temperaturgrenzen der Stellwiderstände zu berücksichtigen.

Für die zwölf dreiphasigen Abgangseinheiten werden insgesamt
36 phasenbezogene Stromregelkreise benötigt. Die einzelnen
Stromregelkreise sind jedoch nicht vollständig voneinander unabhängig.
Alle Abgänge werden aus gemeinsamen Kupferschienen gespeist. Eine
Änderung eines Stellwiderstands verändert deshalb die Gesamtbelastung
der Spannungsquelle und kann über deren innere Impedanz geringfügige
Spannungsänderungen hervorrufen. Diese wirken wiederum auf die übrigen
Abgänge zurück.

Die Regelgrößen und Stellgrößen des nachgelagerten Systems können
vereinfacht zu den Vektoren

\begin{equation}
    \mathbf{I}_{\mathrm{MCC}}
    =
    \begin{bmatrix}
        I_{1,\mathrm{L1}} &
        I_{1,\mathrm{L2}} &
        I_{1,\mathrm{L3}} &
        \ldots &
        I_{12,\mathrm{L3}}
    \end{bmatrix}^{\mathrm{T}}
    \label{eq:stromvektor_mcc}
\end{equation}

und

\begin{equation}
    \mathbf{R}_{\mathrm{Stell}}
    =
    \begin{bmatrix}
        R_{\mathrm{Stell},1,\mathrm{L1}} &
        R_{\mathrm{Stell},1,\mathrm{L2}} &
        R_{\mathrm{Stell},1,\mathrm{L3}} &
        \ldots &
        R_{\mathrm{Stell},12,\mathrm{L3}}
    \end{bmatrix}^{\mathrm{T}}
    \label{eq:stellwiderstandsvektor}
\end{equation}

zusammengefasst werden. Da eine Stellgrößenänderung neben dem
zugeordneten Modulstrom auch die gemeinsamen elektrischen Größen
beeinflussen kann, ist die gesamte Prüfanordnung als gekoppeltes
Mehrgrößensystem zu betrachten.

Für den vorgesehenen Prüfablauf werden die Stellwiderstände zunächst
in eine definierte, strombegrenzende Ausgangsposition gefahren.
Anschließend wird die gemeinsame Versorgungsspannung auf den
vorgesehenen Arbeitspunkt eingestellt und durch die vorgelagerte
Spannungsregelung konstant gehalten. Erst danach werden die
phasenbezogenen Stromregelkreise freigegeben. Die eigentliche
Erwärmungsprüfung kann beginnen, sobald sämtliche Modulströme innerhalb
des festgelegten Toleranzbereichs liegen und ein stabiler Betriebspunkt
erreicht ist.

Die vorliegende Arbeit umfasst keine hardwareseitige Realisierung
dieser Regelungsstruktur. Die Regelkreise werden daher im weiteren
Verlauf durch mathematische Modelle und eine Simulation untersucht.
Die Auswahl der Reglerstruktur, ihre Parametrierung, die Abbildung der
Stellgrößenbegrenzungen und die Untersuchung der gegenseitigen
Beeinflussung erfolgen im Rahmen der Modellbildung und
Plausibilisierung.


% -----------------------------------------------------------------------------
% 4.6 Dimensionierung und erforderliche Stellreserven
% -----------------------------------------------------------------------------
\subsection{Dimensionierung und erforderliche Stellreserven}
\label{sec:dimensionierung_stellreserve}

Ziel der Dimensionierung ist die Festlegung eines Wertebereichs, in
dem der jeweilige Modulstrom sicher eingestellt und während der
Erwärmungsprüfung nachgeführt werden kann. Maßgebend sind dabei das
Übersetzungsverhältnis und die innere Impedanz von Trafo~2, der
Widerstandsbereich von \(R_{\mathrm{Stell}}\), die gemeinsame
Versorgungsspannung sowie der für den jeweiligen Abgang erforderliche
Strombereich.

Da im Rahmen dieser Arbeit keine hardwareseitige Realisierung erfolgt,
wird keine abschließende konstruktive Dimensionierung vorgenommen.
Stattdessen werden die für die Funktionsfähigkeit des Konzepts
notwendigen Zusammenhänge und Auslegungsbedingungen hergeleitet. Diese
bilden anschließend die Grundlage für die Parametrierung des
Simulationsmodells.

Die auf die Primärseite bezogene innere Impedanz von Trafo~2 wird durch

\begin{equation}
    \underline{Z}_{\mathrm{T2},1}
    =
    R_{\mathrm{T2},1}
    +
    \mathrm{j}X_{\sigma,\mathrm{T2},1}
    \label{eq:trafo2_innenimpedanz}
\end{equation}

beschrieben. Darin umfasst \(R_{\mathrm{T2},1}\) die auf die
Primärseite bezogenen ohmschen Wicklungsverluste und
\(X_{\sigma,\mathrm{T2},1}\) die zugehörige Streureaktanz. Zusammen mit
dem transformierten Stellwiderstand ergibt sich die für den
Primärstrom wirksame Gesamtimpedanz zu

\begin{equation}
    \underline{Z}_{\mathrm{ges,T2}}
    =
    R_{\mathrm{T2},1}
    +
    u_{\mathrm{T2}}^{\,2}R_{\mathrm{Stell}}
    +
    \mathrm{j}X_{\sigma,\mathrm{T2},1}.
    \label{eq:gesamtimpedanz_trafo2_dimensionierung}
\end{equation}

Der Betrag des Primärstroms folgt damit aus

\begin{equation}
    I_{\mathrm{T2},1}
    =
    \frac{U_1}
    {
        \sqrt{
            \left(
                R_{\mathrm{T2},1}
                +
                u_{\mathrm{T2}}^{\,2}R_{\mathrm{Stell}}
            \right)^2
            +
            X_{\sigma,\mathrm{T2},1}^{\,2}
        }
    }.
    \label{eq:primaerstrom_trafo2_dimensionierung}
\end{equation}

Gleichung~\eqref{eq:primaerstrom_trafo2_dimensionierung} zeigt, dass
der maximal erreichbare Primärstrom nicht allein durch
\(R_{\mathrm{Stell},\min}\) bestimmt wird. Auch die innere Impedanz von
Trafo~2 begrenzt den Strom. Selbst für einen idealisierten
sekundärseitigen Kurzschluss mit

\begin{equation}
    R_{\mathrm{Stell}} = 0
\end{equation}

verbleibt die Kurzschlussimpedanz des Transformators im Stromkreis. Der
maximal erreichbare Primärstrom beträgt daher

\begin{equation}
    I_{\mathrm{T2},1,\max}
    =
    \frac{U_1}
    {
        \sqrt{
            R_{\mathrm{T2},1}^{\,2}
            +
            X_{\sigma,\mathrm{T2},1}^{\,2}
        }
    }.
    \label{eq:maximalstrom_trafo2}
\end{equation}

Für die spätere Auslegung von Trafo~2 sind somit insbesondere dessen
Kurzschlussimpedanz und relative Kurzschlussspannung erforderlich. Bei
einem einphasig betrachteten Transformator kann die auf die
Primärseite bezogene Kurzschlussimpedanz aus der relativen
Kurzschlussspannung näherungsweise durch

\begin{equation}
    Z_{\mathrm{k},1}
    =
    \frac{u_{\mathrm{k}}}{100}
    \frac{U_{1,\mathrm{N}}}{I_{1,\mathrm{N}}}
    =
    \frac{u_{\mathrm{k}}}{100}
    \frac{U_{1,\mathrm{N}}^{\,2}}{S_{\mathrm{N}}}
    \label{eq:kurzschlussimpedanz_trafo2}
\end{equation}

bestimmt werden. Dabei bezeichnet \(u_{\mathrm{k}}\) die relative
Kurzschlussspannung in Prozent. Ohne diese Kenngröße kann nicht
abschließend beurteilt werden, ob der gewünschte maximale Primärstrom
mit dem vorgesehenen Transformator erreicht wird.

Wird Gleichung~\eqref{eq:primaerstrom_trafo2_dimensionierung} nach dem
Stellwiderstand aufgelöst, ergibt sich der für einen vorgegebenen
Primärstrom erforderliche sekundärseitige Widerstand zu

\begin{equation}
    R_{\mathrm{Stell}}
    =
    \frac{1}{u_{\mathrm{T2}}^{\,2}}
    \left[
        \sqrt{
            \left(
                \frac{U_1}{I_{\mathrm{T2},1}}
            \right)^2
            -
            X_{\sigma,\mathrm{T2},1}^{\,2}
        }
        -
        R_{\mathrm{T2},1}
    \right].
    \label{eq:rstell_aus_primaerstrom}
\end{equation}

Eine reelle Lösung ist nur möglich, wenn

\begin{equation}
    \left(
        \frac{U_1}{I_{\mathrm{T2},1}}
    \right)^2
    \geq
    X_{\sigma,\mathrm{T2},1}^{\,2}
    \label{eq:realisierbarkeitsbedingung_rstell}
\end{equation}

gilt und der nach
Gleichung~\eqref{eq:rstell_aus_primaerstrom} berechnete
Widerstand innerhalb des technisch verfügbaren Stellbereichs liegt.

Für die Abgänge mit Bemessungsströmen von \SI{250}{\ampere},
\SI{400}{\ampere} und \SI{630}{\ampere} übernimmt der vorhandene
V2A-Widerstand den überwiegenden Anteil des Modulstroms. Trafo~2 muss
lediglich die Differenz zwischen dem Strom durch den V2A-Pfad und dem
vorgegebenen Modulstrom bereitstellen. Unter der Annahme eines
überwiegend ohmschen Verhaltens kann der erforderliche Strom des
Zusatzpfads näherungsweise durch

\begin{equation}
    I_{\mathrm{T2},1,\mathrm{erf}}(\vartheta)
    =
    I_{\mathrm{soll}}
    -
    I_{\mathrm{V2A}}(\vartheta)
    \label{eq:erforderlicher_zusatzstrom_temperatur}
\end{equation}

bestimmt werden.

Mit zunehmender Temperatur steigt der elektrische Widerstand des
V2A-Materials. Bei näherungsweise konstanter Phasenspannung sinkt
dadurch der Strom \(I_{\mathrm{V2A}}(\vartheta)\), während der
erforderliche Zusatzstrom nach
Gleichung~\eqref{eq:erforderlicher_zusatzstrom_temperatur} zunimmt. Die
notwendige Nachführreserve ergibt sich folglich aus der Differenz
zwischen dem V2A-Strom im kalten und im maximal betrachteten warmen
Zustand:

\begin{equation}
    \Delta I_{\mathrm{th}}
    =
    I_{\mathrm{V2A}}(\vartheta_{\min})
    -
    I_{\mathrm{V2A}}(\vartheta_{\max}).
    \label{eq:thermische_stromreserve}
\end{equation}

Zur thermisch bedingten Stromänderung ist eine zusätzliche Reserve für
Modellabweichungen, Bauteiltoleranzen und Schwankungen des
Arbeitspunkts vorzusehen. Der maximal erforderliche Strom des
Zusatzpfads wird deshalb durch

\begin{equation}
    I_{\mathrm{T2},1,\mathrm{aus}}
    =
    I_{\mathrm{T2},1,\mathrm{erf,max}}
    +
    \Delta I_{\mathrm{Reserve}}
    \label{eq:auslegungsstrom_zusatzpfad}
\end{equation}

festgelegt. Der Wert \(\Delta I_{\mathrm{Reserve}}\) ist im weiteren
Verlauf anhand der Simulationsergebnisse zu bestimmen.

Da der transformatorgestützte Zusatzpfad ausschließlich einen
positiven Strombeitrag liefern kann, muss bereits im kalten Zustand
sichergestellt werden, dass der Gesamtstrom nicht oberhalb des
zulässigen Bereichs liegt. Die entsprechende Bedingung wurde in
Gleichung~\eqref{eq:bedingung_kaltzustand_grosse_abgaenge}
formuliert. Im maximal betrachteten warmen Zustand muss dagegen die in
Gleichung~\eqref{eq:bedingung_warmzustand_grosse_abgaenge}
angegebene Bedingung erfüllt sein. Der V2A-Widerstand ist somit so
vorzudimensionieren, dass Trafo~2 im kalten Zustand nicht bereits an
seiner unteren Stellgrenze und im warmen Zustand nicht an seiner
oberen Stellgrenze betrieben wird.

Für einen gewählten Arbeitspunkt
\(R_{\mathrm{Stell},0}\) können die verfügbaren
Widerstandsreserven durch

\begin{align}
    \Delta R_{\mathrm{Stell},\downarrow}
    &=
    R_{\mathrm{Stell},0}
    -
    R_{\mathrm{Stell},\min},
    \label{eq:rstell_reserve_unten}
    \\[1ex]
    \Delta R_{\mathrm{Stell},\uparrow}
    &=
    R_{\mathrm{Stell},\max}
    -
    R_{\mathrm{Stell},0}
    \label{eq:rstell_reserve_oben}
\end{align}

beschrieben werden. Die Reserve in Richtung kleinerer Widerstände
ermöglicht eine Erhöhung des Modulstroms. Die Reserve in Richtung
größerer Widerstände ermöglicht dessen Verringerung.

Bei den Abgängen mit Bemessungsströmen von \SI{40}{\ampere} und
\SI{70}{\ampere} ist kein paralleler V2A-Widerstand vorgesehen. Trafo~2
muss daher den vollständigen Modulstrom übertragen. Für diese Abgänge
ist nicht nur eine thermisch bedingte Stromdifferenz, sondern ein
vollständiger Stromstellbereich bereitzustellen. Der Bemessungsstrom
muss einschließlich der erforderlichen oberen und unteren
Stellreserve innerhalb des erreichbaren Bereichs liegen:

\begin{equation}
    I_{\mathrm{T2},1,\min}
    \leq
    I_{\mathrm{N}}
    -
    \Delta I_{\mathrm{Reserve},-}
    <
    I_{\mathrm{N}}
    <
    I_{\mathrm{N}}
    +
    \Delta I_{\mathrm{Reserve},+}
    \leq
    I_{\mathrm{T2},1,\max}.
    \label{eq:stromreserve_kleine_abgaenge}
\end{equation}

Die untere Stromreserve ermöglicht eine Verringerung des Stroms bei
einem zu hohen Ausgangswert. Die obere Stromreserve wird benötigt, um
Widerstandserhöhungen der Wicklungen, Kontakte und Verbindungen sowie
Abweichungen der Versorgungsspannung ausgleichen zu können.

Da alle Abgänge an einer gemeinsamen Phasenspannung betrieben werden,
müssen die Übersetzungsverhältnisse und Widerstandsbereiche so gewählt
werden, dass die jeweiligen Arbeitspunkte gleichzeitig erreichbar
sind. Für jeden Kanal muss bei der vorgesehenen Phasenspannung gelten:

\begin{equation}
    I_{k,\varphi,\min}
    \leq
    I_{\mathrm{soll},k,\varphi}
    \leq
    I_{k,\varphi,\max}.
    \label{eq:gemeinsamer_spannungsarbeitspunkt}
\end{equation}

Kann diese Bedingung nicht mit einer einheitlichen Ausführung erfüllt
werden, sind für die unterschiedlichen Modulströme angepasste
Übersetzungsverhältnisse oder unterschiedliche Stellwiderstandsbereiche
vorzusehen. Eine identische Ausführung von Trafo~2 und
\(R_{\mathrm{Stell}}\) für sämtliche Stromklassen wird daher nicht
vorausgesetzt.

Neben dem Widerstandsbereich ist die Dauerstrom- und
Dauerleistungsfähigkeit der Komponenten zu berücksichtigen. Für das
verwendete Übersetzungsverhältnis

\begin{equation}
    u_{\mathrm{T2}}
    =
    \frac{U_1}{U_2}
    =
    \frac{N_1}{N_2}
\end{equation}

gelten im idealisierten Fall

\begin{align}
    U_2
    &=
    \frac{U_1}{u_{\mathrm{T2}}},
    \label{eq:sekundaerspannung_trafo2}
    \\
    I_2
    &=
    u_{\mathrm{T2}} I_1.
    \label{eq:sekundaerstrom_trafo2}
\end{align}

Durch eine sekundärseitige Spannungserhöhung wird der dort fließende
Strom gegenüber dem hohen Primärstrom reduziert. Dadurch kann der
Stellwiderstand mit höheren Widerstandswerten und geringeren Strömen
ausgeführt werden. Seine Verlustleistung beträgt

\begin{equation}
    P_{\mathrm{Stell}}
    =
    I_2^{\,2}R_{\mathrm{Stell}}.
    \label{eq:verlustleistung_stellwiderstand_dimensionierung}
\end{equation}

Da sowohl \(I_2\) als auch \(R_{\mathrm{Stell}}\) vom eingestellten
Arbeitspunkt abhängen, muss die maximale Verlustleistung über den
gesamten Stellbereich ermittelt werden:

\begin{equation}
    P_{\mathrm{Stell},\max}
    =
    \max_{
        R_{\mathrm{Stell},\min}
        \leq
        R_{\mathrm{Stell}}
        \leq
        R_{\mathrm{Stell},\max}
    }
    \left\{
        I_2^{\,2}(R_{\mathrm{Stell}})
        R_{\mathrm{Stell}}
    \right\}.
    \label{eq:maximale_verlustleistung_stellwiderstand}
\end{equation}

Eine alleinige Betrachtung der beiden Endlagen des Stellwiderstands
ist nicht grundsätzlich ausreichend, da die maximale Verlustleistung
auch innerhalb des Stellbereichs auftreten kann.

Trafo~2 und \(R_{\mathrm{Stell}}\) sind für den Dauerbetrieb während
der gesamten Erwärmungsprüfung auszulegen. Für die kleinen Abgänge ist
die Primärwicklung von Trafo~2 für den vollständigen Modulstrom
einschließlich Stellreserve zu bemessen. Bei den großen Abgängen ist
dagegen der maximal erforderliche Zusatzstrom maßgebend. Zusätzlich
müssen die zulässige magnetische Flussdichte, die Wicklungserwärmung,
die Isolationsbeanspruchung und die thermische Rückwirkung der
Stellwiderstände auf den eigentlichen Prüfling berücksichtigt werden.

Für eine abschließende Dimensionierung werden mindestens folgende
Kenngrößen benötigt:

\begin{itemize}
    \item Bemessungsleistung und Übersetzungsverhältnis von Trafo~2,
    \item relative Kurzschlussspannung beziehungsweise
          Kurzschlussimpedanz von Trafo~2,
    \item ohmsche Wicklungswiderstände und Streureaktanzen,
    \item minimaler und maximaler Wert von \(R_{\mathrm{Stell}}\),
    \item zulässiger Dauerstrom und zulässige Verlustleistung des
          Stellwiderstands,
    \item maximal betrachtete Temperatur der V2A-Widerstände,
    \item Widerstände der MCC-Module, Kontakte und Verbindungen,
    \item vorgesehene Versorgungsspannung und zulässige
          Stromabweichung.
\end{itemize}

Da die konkreten Kurzschlussdaten von Trafo~2 gegenwärtig nicht
vollständig vorliegen, werden die entsprechenden Parameter im
Simulationsmodell zunächst als Modellannahmen festgelegt und
transparent ausgewiesen. Die Simulation dient damit der Untersuchung
der prinzipiellen Funktionsfähigkeit und der erforderlichen
Stellreserven. Sie stellt keine abschließende Bau- oder
Komponentenfreigabe für eine spätere hardwareseitige Umsetzung dar.




Zusammenfassend ergeben sich aus der konzeptionellen Untersuchung die
Anforderungen an den Übersetzungsfaktor und die innere Impedanz von
Trafo~2, den Stellbereich und die Belastbarkeit von
\(R_{\mathrm{Stell}}\) sowie an die erforderlichen Strom- und
Positionsreserven. Da der entwickelte Aufbau im Rahmen dieser Arbeit
nicht als reale Prüfanordnung umgesetzt wird, erfolgt der Nachweis der
grundsätzlichen Funktionsfähigkeit anhand eines mathematischen
Simulationsmodells. Die hierzu erforderliche Modellbildung ist
Gegenstand des folgenden Kapitels.
